

% This section will go away upon submission
\section*{Reviewers Comments}

\begin{markdown}

## Reviewer 1

- [x] Software team itself: team size, composition, and collaboration practices. **in Background section**

## Reviewer 2

- [x] Figure 2 is challenging to decipher in its regular print format; widening the bars **Changed by Philip**

- [x] Team interactions (see reviewer 1 comment) **Added paragraph at the end of the Background section**

## Reviewer 3 

- [x] Comment on maintainer and contributor experience. **in Background section, but not sure how to capture subjective information. Perhaps for the presentation.**

- [x] For non-practitioners, should say what the difference between =qmcpack= and =qmcpack_complex= is, or better /why/ it's important.  It's implicitly mentioned in "Compile-time to runtime variants". **added real/complex number calculations**

- [x] Does the CI time of ~1 hour pose any barriers to development? **Probably for the presentation? Short answer is No.**

- [x] How different is the Docker environment for CI from (core) developer's environments? **Probably for the presentation? Not sure how it is relevant since the target is HPC**

- [x] Formatting of "std :: unique ptr and std :: shared ptr" **Changed in the text in Section III.C**

- [x] Why is code coverage sensitive to the gcov version? **Good question, we didn't investigate further. Perhaps bugs?**

- [x] How do both regular and infrequent contributors feel about the changes to CI? **Presentation? This is highly subjective**

- [x] Has the code review process changed at all as part of CI modernization? **Presentation? The code review process is somehow orthogonal?**

- [x] In terms of the latest Spack+Docker+GHA infrastructure, where is the largest maintenance cost?  Is there any new overhead in maintaining this infrastructure?  Being proactive by preventing outside world changes from introducing technical debt means that there /is/ continued maintenance required rather than "one-and-done". **Not sure if I follow the last sentence. The biggest overhead is keeping up with LLVM versions in GPU CI, CPU CI is a very virtual reproducible environment in which we don't have to manage the hardware and software in the short-term. Perhaps for the presentation?**

## Reviewer 4

- [x] At the start of §2, instead of just the number of citations of the reference paper, also include some relevant examples of applications to demonstrate the breadth and scope of research or work enabled by
QMCPACK. **This is covered in the introduction**
   
- [x] Why is spack the preferred installer? **added in Background section, as Spack targets supercomputers**

- [x] In Table 3, what are the operating systems of Sulfur and Nitrogen? **Added in subsection Strategic CI expansion on GitHub Actions**

- [x] In figure 1, subcategorizing the failure types might aide the discussion of the 2023 Q1 failures more clearly. ** This would create a lot of data handling from GitHub Actions with little added value**

- [x] In Section 3.c, figure 3: why was April of 2023 so good for code coverage? **Added due to increase in GCC version**

- [x] The related work section is wonderful resource and fits nicely in the scope of the conference, but feels out of place with the rest of the manuscript. A lead-in paragraph may be enough to justify it. **An lead-in sentence was added to motivate the introduction of this section**

- [x] Some discussion on scope creep and hardware support would be interesting. What were the main challenges in the CUDA/HIP/ROCm/SYCL development cycles? **Leave for the presentation or extended version?**


\end{markdown}